% ═══════════════════════════════════════════════════════════════════
% §6  Air Source Heat Pump Boiler (Dynamic Model)
% Source: AirSourceHeatPumpBoiler.py
% ═══════════════════════════════════════════════════════════════════
\section{Air source heat pump boiler}
\label{sec:ashpb}

This section describes the dynamic ASHPB model implemented in \texttt{AirSourceHeatPumpBoiler.py}. It integrates a vapour-compression refrigerant cycle, physics-based heat exchangers, a Differential Evolution optimiser, and a storage tank with hysteresis control.

\subsection{Assumptions}
\begin{enumerate}
  \item The refrigerant cycle operates in quasi-steady state within each time step.
  \item Refrigerant properties are computed via CoolProp (real-fluid EOS).
  \item Superheat $\Delta T_\text{sh}$ and subcooling $\Delta T_\text{sc}$ are constant.
  \item The compressor follows a displacement-based mass flow model with isentropic efficiency.
  \item The condenser is a submerged coil (\emph{UA}$\cdot\Delta T$ model, not LMTD).
  \item The evaporator uses an air-side energy balance (cross-flow).
\end{enumerate}

\subsection{Refrigerant cycle state points}

Saturation temperatures are set by the source/sink conditions:
\begin{align}
  T_\text{eva,sat} &= T_0 - \Delta T_\text{ref,evap}  \label{eq:ashp-Teva}\\
  T_\text{cond,sat} &= T_\text{w,tank} + \Delta T_\text{ref,cond} \label{eq:ashp-Tcond}
\end{align}
where $\Delta T_\text{ref,evap}$ and $\Delta T_\text{ref,cond}$ are the Differential Evolution optimisation variables.

Compressor inlet (superheated):
\begin{equation}\label{eq:ashp-cmp-in}
  T_\text{cmp,in} = T_\text{eva,sat} + \Delta T_\text{sh}
  \,,\quad
  h_\text{cmp,in} = h(T_\text{cmp,in},\,P_\text{eva})
\end{equation}

Compressor outlet (isentropic + efficiency):
\begin{equation}\label{eq:ashp-cmp-out}
  h_\text{cmp,out} = h_\text{cmp,in} + \frac{h_\text{cmp,out,s} - h_\text{cmp,in}}{\eta_\text{isen}}
\end{equation}

\subsection{Heat exchanger models}

\paragraph{Condenser (submerged coil).}
\begin{equation}\label{eq:ashp-Qcond}
  Q_\text{cond} = UA_\text{cond}\,\Delta T_\text{ref,cond}
\end{equation}

\paragraph{Evaporator (air-side energy balance).}
\begin{equation}\label{eq:ashp-Qevap}
  Q_\text{evap} = \rho_a\,c_a\,\dot{V}_\text{fan}\,(T_\text{a,in} - T_\text{a,out})
\end{equation}
The fan flow rate $\dot{V}_\text{fan}$ is determined by matching the air-side capacity to the refrigerant-side evaporation duty $Q_\text{ref,evap} = \dot{m}_\text{ref}(h_\text{eva,out} - h_\text{exp,out})$.

\subsection{Compressor model}

Mass flow rate from displacement volume:
\begin{equation}\label{eq:ashp-mdot}
  \dot{m}_\text{ref} = \rho_\text{cmp,in}\,V_\text{disp}\,\frac{n_\text{cmp}}{60}
\end{equation}
where $n_\text{cmp}$ [RPM] is determined by the condenser heat balance.

Power use:
\begin{equation}\label{eq:ashp-Ecmp}
  E_\text{cmp} = \dot{m}_\text{ref}\,(h_\text{cmp,out} - h_\text{cmp,in})
\end{equation}

\subsection{Differential Evolution optimisation}

\paragraph{Objective.}
Minimise total power use:
\begin{equation}\label{eq:ashp-obj}
  \min_{\Delta T_\text{ref,cond},\,\Delta T_\text{ref,evap}} \;
  E_\text{tot} = E_\text{cmp} + E_\text{fan,ou}
\end{equation}

\paragraph{Constraints.}
\begin{itemize}
  \item Heat balance: $Q_\text{cond} = \dot{m}_\text{ref}(h_\text{cmp,out} - h_\text{exp,in})$
  \item Minimum/maximum fan flow: $0.1\,\dot{V}_\text{design} \le \dot{V}_\text{fan} \le \dot{V}_\text{design}$
  \item Compressor speed bounds: $n_\text{min} \le n_\text{cmp} \le n_\text{max}$
\end{itemize}

\subsection{Storage tank energy balance}
\label{sec:ashpb-tank}

The storage tank is modelled as a lumped capacitance with variable mass.
The governing ODE is:
\begin{equation}\label{eq:ashp-tank-ode}
  C \frac{dT}{dt} = \dot{Q}_\text{gain} - UA_\text{tank}\,(T - T_0)
\end{equation}
where $C = c_w \rho_w V_\text{tank} \cdot \ell$ is the effective thermal capacitance
(scaled by the tank level $\ell \in (0,1]$),
and $\dot{Q}_\text{gain}$ aggregates all source terms
(condenser, UV lamp, STC, refill enthalpy).

\paragraph{Fully implicit discretisation.}
All temperature-dependent terms are evaluated at $n{+}1$, and the
coupled mass--energy system is solved simultaneously via Newton--Raphson
(\texttt{scipy.optimize.fsolve}).

\textbf{Energy residual:}
\begin{equation}\label{eq:ashp-implicit-energy}
  r_E = C_\text{full}\,\ell^{n+1}\,T^{n+1}
      - C_\text{full}\,\ell^{n}\,T^{n}
      - \Delta t\,\bigl[
          \dot{Q}_\text{HP} + E_\text{UV} + E_\text{pump}
          + \dot{Q}_\text{STC}(T^{n+1})
          + \dot{Q}_\text{refill}(T^{n+1})
          - UA\,(T^{n+1} - T_0)
        \bigr]
\end{equation}
where
$\dot{Q}_\text{refill}(T) = c_w \rho_w \dot{V}_\text{in}(T_\text{refill} - T)$
and
$\dot{Q}_\text{STC}(T)$ is the net STC gain evaluated at tank temperature~$T$.

\textbf{Mass residual:}
\begin{equation}\label{eq:ashp-implicit-mass}
  r_M = \ell^{n+1} - \ell^{n}
      - \frac{\dot{V}_\text{in} - \dot{V}_\text{out}(T^{n+1})}{V_\text{full}}\,\Delta t
\end{equation}
where $\dot{V}_\text{out}(T) = \alpha(T)\,\dot{V}_\text{mix}$
with the mixing ratio
$\alpha(T) = \text{clip}\bigl(\frac{T_\text{mix}-T_\text{in}}{T - T_\text{in}},\;0,\;1\bigr)$.

The unknowns $[T^{n+1},\,\ell^{n+1}]$ are found by solving
$\mathbf{r} = \mathbf{0}$ using \texttt{fsolve}.
This scheme is unconditionally stable and handles strong coupling
between variable tank mass and temperature.

\subsection{Exergy analysis}

Subsystem exergy balances follow the input--destruction--output pattern~\cite{tsatsaronis2007}:

\paragraph{Refrigerant loop.}
\begin{equation}\label{eq:ashp-Xr}
  X_\text{cmp} - X_\text{c,r} = X_\text{r,iu} + X_\text{r,ou}
\end{equation}

\paragraph{Indoor unit (condenser side).}
\begin{equation}\label{eq:ashp-Xiu}
  X_\text{r,iu} - X_\text{c,iu} = X_\text{w,out} - X_\text{w,in}
\end{equation}

\paragraph{Overall exergy efficiency.}
\begin{equation}\label{eq:ashp-etaX}
  \eta_{X,\text{sys}} = \frac{X_\text{w,out} - X_\text{w,in}}{X_\text{cmp} + X_\text{fan,iu} + X_\text{fan,ou}}
\end{equation}
